\documentclass[a4paper, oneside]{article}
\usepackage{graphicx}
\usepackage{amsthm}
\usepackage{amsmath}
\usepackage{amssymb}
\usepackage[a4paper,
            bindingoffset=0.2in,
            left=2cm,
            right=2cm,
            top=2cm,
            bottom=2cm,
            footskip=.25in]{geometry}
\usepackage[italian]{babel}
\usepackage{pgfplots}
\usepackage{tabularx}
\usepackage{tikz}
\usepackage{wrapfig}
\usepackage{color}
\usepackage[d]{esvect}
\definecolor{page}{rgb}{0.129,0.157,0.212}
\pagecolor{page}
\color{white}
\graphicspath{ {./images/} }
\usetikzlibrary{shapes.geometric}
\usetikzlibrary{datavisualization}
\usetikzlibrary{datavisualization.formats.functions}
\usetikzlibrary{patterns}
\pgfplotsset{width=10cm,compat=1.18}

\title{Esercizi di Fisica II}
\author{Tommaso Miliani}
\date{06-03-26}

\begin{document}
\newtheoremstyle{theoremEnv}
                {}          % Space above
                {}          % Space below
                {\slshape}  % Body font
                {}          % Indent amount
                {\bfseries} % Head font
                {.}         % Punctuation after head
                {\newline}  % Space after theorem head
                {}          % Theorem head spec
\theoremstyle{theoremEnv}

\newtheorem{definition}{Definizione}[section]
\newtheorem{theorem}{Teorema}[section]
\newtheorem{lemma}{Proposizione}[section]
\newtheorem{observation}{Osservazione}[section]
\newtheorem{corollary}{Corollario}[theorem]
\newtheorem{example}{Esempio}[section]
\newtheorem{remark}{Enunciato}[section]

\maketitle

\section{Esercizi}
\begin{example}[20/02/2012]
    Se si volesse calcolare il potenziale dentro la sfera, è dato 
    \begin{gather*}
        z = R\cos\theta \quad dQ = \sigma(\theta)dl Rd\theta = \sigma_0\cos\theta dl d\theta R\\
        r(\theta) = R\sin\theta \quad 
    \end{gather*}
    Dunque 
    \begin{gather*}
        V = \int_{}^{} dl \int_{}^{} d\theta \frac{dQ}{4\pi\epsilon_0 R} = \frac{1}{4\pi \epsilon_0} \int_{}^{} dl \int_{0}^{\pi} d\theta \sigma_0\cos\theta = \frac{1}{4\pi\epsilon_0} \int_{0}^{\pi} d\theta \sigma_0 \cos\theta \pi R^{2} \sin\theta = 0 
    \end{gather*}
    Dove $dl$ è un tratto infinitesimo sulla sfera: l'integrale su $dl$ è dunque l'integrale lungo 
    la spira mentre quello in $d\theta$ è su tutta la sfera. Il potenziale dentro la sfera è 
    dunque zero. Il campo al centro della sfera si trova in maniera simile: sfruttando le stesse 
    considerazioni geometriche si può trovare 
    infatti il contributo infinitesimo di campo 
    \begin{gather*}
        dE_z = \frac{dQ}{4\pi\epsilon_0 R^{2}}\cos\theta
    \end{gather*}
    Il campo totale $\vv{E}$ sarà dunque solo lungo $\hat{z}$:
    \begin{gather*}
        E_z = \int_{0}^{\pi} \int_{SPIRA} \frac{\sigma_0 \cos\theta R d\theta dl}{4\pi\epsilon_0 R^{2}} \cos\theta = \frac{\sigma_0}{4\pi\epsilon_0 R} \int_{0}^{\pi} d\theta \cos^{2}\theta \int_{SPIRA}^{} dl      
    \end{gather*}  
    Ottenendo 
    \begin{gather*}
        E_z = \int_{0}^{\pi} d\theta \cos^{2}\theta 2\pi R\sin\theta = \frac{\sigma_0}{3\epsilon_0}
    \end{gather*}
\end{example}

Considerata una carica negativa ed una positiva distanziate da una stessa quantità $d$. Si definisce il dipolo elettrico 
$\vv{P} = q\vv{d}$, ossia il vettore spostamento tra le cariche. Il campo elettrico di un dipolo è importante perché è il campo che 
genera un atomo. Si vuole ora colcolare il campo elettrico totale di un sistema di questa tipologia e gli effetti meccanici che
agiscono su di esso. Il campo elettrico in un punto generico $r >> d$, 
\begin{gather*}
    \vv{r} = x\hat{x} + y\hat{y} + z\hat{z} \qquad \vv{r_+} = \frac{d}{2} \hat{z} \qquad \vv{r_-} = - \frac{d}{2}\hat{z}        
\end{gather*}
Si vuole ora determinare il potenziale, il quale sarà 
\begin{gather*}
    V = V_+ + V_- = \frac{q}{4\pi\epsilon_0} \left(\frac{1}{\left| \vv{r} - \vv{r_+}   \right| } - \frac{1}{\left| \vv{r} - \vv{r_-}   \right| }\right)
\end{gather*}
I vettori 
\begin{align*}
    \vv{r} - \vv{r_+} &= x\hat{x} + y\hat{y} + \left(z - \frac{d}{2}\right)    \hat{z}  &\left| \vv{r} - \vv{r_+}   \right| = \sqrt{ x^{2} + y^{2} + \left(z - \frac{d}{2}\right)^{2}} \\
    \vv{r} - \vv{r_-} &= x\hat{x} + y\hat{y} + \left(z + \frac{d}{2}\right) \hat{z}      &\left| \vv{r} - \vv{r_-}  \right| = \sqrt{x^{2} + y^{2} + \left(z + \frac{d}{2}\right)^{2}} 
\end{align*}
Dunque
\begin{gather*}
    \frac{1}{\left| \vv{r} - \vv{r_+}   \right| } = \frac{1}{\sqrt{x^{2} + y^{2} + z^{2} - zd + \frac{d^{2}}{4}} } = \frac{1}{\sqrt{r^{2} \left(1 + \frac{d^{2}}{4r^{2}} - \frac{zd}{r^{2}}\right) }} =  \frac{1}{r} \frac{1}{\left(1 + \frac{d^{2}}{4r^{2}} - \frac{zd}{r^{2}}\right)} \\
    \frac{1}{\left| \vv{r} - \vv{r_-}   \right| } = \frac{1}{\sqrt{x^{2} + y^{2} + z^{2} + zd + \frac{d^{2}}{4}} } = \frac{1}{\sqrt{r^{2} \left(1 + \frac{d^{2}}{4r^{2}} + \frac{zd}{r^{2}}\right) }} =  \frac{1}{r} \frac{1}{\left(1 + \frac{d^{2}}{4r^{2}} + \frac{zd}{r^{2}}\right)} 
\end{gather*}
SI può ora approssimare che $d^{2} << r^{2}$, ossia secondo l'approssimazione dipolo. Si sviluppa dunque 
con Taylor considerando $\frac{zd}{r^{2}}$ come $\epsilon$, si può ignorare il termine $d^{2}$ e sviluppare:
\begin{gather*}
    \frac{1}{\sqrt{1 \pm \epsilon} } \approx 1 \mp \frac{1}{2}\epsilon 
\end{gather*}
Dunque
\begin{gather*}
    V = \frac{q}{4\pi\epsilon_0 } \left(\frac{1}{2} \left(1 + \frac{1}{2} \epsilon\right) - \frac{1}{2}\left(1 - \frac{1}{2}\epsilon\right)\right) = \frac{q}{4\pi\epsilon_0} \frac{\epsilon}{r}
\end{gather*}
Allora
\begin{gather*}
    V = \frac{q}{4\pi\epsilon_0} \frac{1}{r} \frac{zd}{r^{2}}
\end{gather*}
Il vettore 
\begin{gather*}
    \left| \vv{P} \right| = qd \quad z = r\cos\theta \ \Longrightarrow \ zqd = \left| \vv{P}  \right| \left| \vv{r}  \right| \cos\theta = \vv{P} \cdot \vv{r}     
\end{gather*}
Allora
\begin{gather*}
    V = \frac{\vv{P} \cdot \vv{r}  }{4\pi\epsilon_0 r^{3}}
\end{gather*}
Ossia il potenziale elettrico generato da un dipolo elettrico in un generico punto dello spazio. Avrà dunque un campo 
non nullo:
\begin{gather*}
    \vv{E} = -\vv{\nabla}V = \frac{\partial V}{\partial r}\hat{r} + \frac{1}{r} \frac{\partial V}{\partial \theta} \hat{\theta} + \frac{1}{r\sin\theta}\frac{\partial V}{\partial \phi}\hat{\phi}        
\end{gather*}
Dove l'angolo $\phi $è l'angolo che la proiezione del vettore $r$ sul piano $xy$ ha rispetto all'asse $x$. Non si ha una dipendenza da $\phi$
in quanto non cambia l'orientazione del vettore. Tuttavia il campo dipenderà anche dall'angolo $\theta$:
\begin{align*}
    E_r &= -\frac{\partial V}{\partial r} = \frac{\left| P \right| \cos\theta }{2\pi\epsilon_0 r^{3}} \\
    E_\theta &= -\frac{1}{r} \frac{\partial V}{\partial \theta} = -\frac{1}{r} \frac{\left| P \right| (-\sin\theta) }{ 4\pi\epsilon_0 r^{2}} = \frac{\left| P \right| }{4\pi\epsilon_0 r^{3}} \sin\theta
\end{align*}
Allora il campo totale
\begin{gather*}
    \vv{E}  = \frac{\left| P \right| \cos\theta }{2\pi\epsilon_0 r^{3}} \hat{r}  + \frac{\left| P \right| }{4\pi\epsilon_0 r^{3}} \sin\theta \hat{\theta} 
\end{gather*}
Si può traformare il vettore $\hat{r}$:
\begin{gather*}
    \hat{r} = \begin{pmatrix}
        \cos\theta & \sin\theta \\
        -\sin\theta & \cos\theta
    \end{pmatrix}  \begin{pmatrix}
        0 \\
        1
    \end{pmatrix} = \sin\theta \hat{i} + \cos\theta \hat{k} 
\end{gather*} 
E il versore $\hat{\theta}$:
\begin{gather*}
    \hat{\theta} = \begin{pmatrix}
        \cos\theta & \sin\theta \\
        -\sin\theta & \cos\theta 
    \end{pmatrix} \begin{pmatrix}
        1 \\
        0 
    \end{pmatrix} = \cos\theta \hat{i} - \sin\theta \hat{k}  
\end{gather*} 
Dunque 
\begin{align*}
    \cos\theta \hat{r} &= \sin\theta \cos\theta \hat{i} + \cos^{2}\theta \hat{k} \\
    \sin\theta \hat{\theta} &= \cos\theta \sin\theta \hat{i} - \sin^{2} \theta \hat{k}     \\
    2\cos\theta \hat{r} + \sin\theta \hat{\theta} &= 3\cos\theta \hat{r} - \cos\theta \hat{r} + \sin\theta \hat{\theta} = 3\cos\theta \hat{r} - \hat{k}        
\end{align*}
Questo ci interessa perché stiamo cercando delle coordinate che lungo $r$ e lungo $k$, sono cilindriche 
nella direzione del dipolo:
\begin{gather*}
    \vv{E} = \frac{\left| P \right| }{4\pi\epsilon_0 r^{3}} \left(3 \cos\theta \hat{r} - \hat{k}  \right) = \frac{1}{4\pi\epsilon_0} \left(3 \frac{(\vv{P} \cdot \vv{r}  ) \vv{r} }{r^{5}} - \frac{\vv{P} }{r^{3}}\right) 
\end{gather*}
Questo campo prende il nome di \textbf{campo di dipolo}.

\subsection{Effetti meccanici di un dipolo}
Gli effetti meccanici di un dipolo sono i seguenti
\begin{gather*}
    \vv{F^{-}} = -q\vv{E(\vv{r} )}   \qquad \vv{F^{+}} = q\vv{E}(\vv{r} + \vv{d}  )  
\end{gather*}
Si può anche determinare il momento che fa il dipolo
\begin{gather*}
    \vv{m} = \vv{r} \times \vv{F^{-}} + (\vv{r} + \vv{d})  \times \vv{F^{+}} = q\left((\vv{r} + \vv{d}  ) \times \vv{E}(\vv{r} + \vv{d}  ) - \vv{r} \times \vv{E}(\vv{r} )  \right)
\end{gather*}
Le componenti del campo sono 
\begin{gather*}
    E_x(\vv{r} + \vv{d}  ) = \vv{E_x} (\vv{r} ) + \frac{\partial E_x}{\partial x} dx + \frac{\partial E_x}{\partial y} dy + \frac{\partial E_x}{\partial z} dz   
\end{gather*}
Dunque
\begin{gather*}
    \vv{E}(\vv{r} + \vv{d}  ) = \vv{E}(\vv{r} ) + (\vv{d} \cdot \vv{\nabla}  )  \vv{E} 
\end{gather*}
Sviluppando al primo ordine la forza totale è nulla, ma il momento di dipolo 
non è nullo ed è l'analogo con il campo elettrico del momento delle forze in meccanica:
\begin{gather*}
    \left\{\begin{array}{l}
        \vv{F} = (\vv{P} \cdot \vv{\nabla}  ) \vv{E} \\
        \vv{m} = \vv{P} \times \vv{E}     
    \end{array}\right.
\end{gather*}

\end{document}